% Paste this complete file into main.tex in Overleaf.
% Compile with pdfLaTeX. No external files or shell escape are required.
\documentclass[a4paper,10pt,twoside,openany]{book}

\usepackage[T1]{fontenc}
\usepackage[utf8]{inputenc}
\usepackage{helvet}
\renewcommand{\familydefault}{\sfdefault}
\usepackage[top=2.2cm,bottom=2.5cm,inner=2.4cm,outer=1.8cm,
  includeheadfoot,headheight=22pt]{geometry}
\usepackage{xcolor,tikz}
\usetikzlibrary{calc,positioning,arrows.meta}
\usepackage[most]{tcolorbox}
\usepackage{fancyhdr,titlesec,titletoc}
\usepackage{amsmath,amssymb,mathtools}
\usepackage{enumitem,pifont,microtype}
\usepackage{listings}
\usepackage[hidelinks]{hyperref}
\setlength{\parindent}{0pt}
\setlength{\parskip}{5pt}
\setlength{\emergencystretch}{2em}

% Colours from the supplied book template.
\definecolor{Cap1Blue}{HTML}{1E9ED7}
\definecolor{Cap2Orange}{HTML}{F07E26}
\definecolor{Cap3Green}{HTML}{5BB85D}
\definecolor{Cap4Purple}{HTML}{7C5CBF}
\definecolor{AccentYellow}{HTML}{FFD93D}
\definecolor{LightBlue}{HTML}{DAEEF8}
\definecolor{LightOrange}{HTML}{FDEBD0}
\definecolor{LightGreen}{HTML}{D5F5D5}
\definecolor{LightPurple}{HTML}{EAE4F7}
\definecolor{DarkText}{HTML}{1A1A2E}
\definecolor{MidGray}{HTML}{6C757D}
\definecolor{BoxGray}{HTML}{F4F6F8}
\definecolor{RuleGray}{HTML}{DEE2E6}
\color{DarkText}
\newcommand{\ChapColor}{Cap1Blue}
\newcommand{\ChapColorLight}{LightBlue}
\newcommand{\setChapColor}[2]{%
  \renewcommand{\ChapColor}{#1}%
  \renewcommand{\ChapColorLight}{#2}}

% Running heads and page numbers.
\fancypagestyle{main}{%
  \fancyhf{}
  \fancyhead[LE]{\small\color{\ChapColor}\nouppercase{\leftmark}}
  \fancyhead[RO]{\small\color{MidGray}\nouppercase{\rightmark}}
  \fancyfoot[LE,RO]{%
    \tikz[baseline=(p.base)]\node[
      fill=\ChapColor,text=white,rounded corners=5pt,
      minimum width=0.65cm,minimum height=0.5cm,
      font=\small\bfseries](p){\thepage};}
  \fancyfoot[C]{\tiny\color{MidGray}FIPpy: ionisation physics}
  \renewcommand{\headrulewidth}{0pt}
  \renewcommand{\footrulewidth}{0pt}}
\pagestyle{main}
\fancypagestyle{plain}{%
  \fancyhf{}
  \fancyfoot[C]{\small\color{MidGray}\thepage}
  \renewcommand{\headrulewidth}{0pt}}

% Chapter and section headings.
\newcommand{\chaptertitleformat}[1]{%
  \begin{tikzpicture}[remember picture,overlay]
    \fill[\ChapColor] (-0.5cm,0) rectangle (0,-4.2cm);
    \fill[\ChapColor] (1.45cm,-1.6cm) circle (1.2cm);
    \fill[white,opacity=0.15] (1.45cm,-1.6cm) circle (0.95cm);
    \node[text=white,font=\fontsize{44}{48}\selectfont\bfseries]
      at (1.45cm,-1.6cm){\thechapter};
    \node[anchor=north west,text=\ChapColor,
      font=\fontsize{25}{29}\selectfont\bfseries,
      align=left,text width=9.5cm] at (3.1cm,-0.4cm){#1};
    \fill[\ChapColor,opacity=0.12]
      (-0.5cm,-3.5cm) rectangle (\textwidth,-4.2cm);
  \end{tikzpicture}
  \vspace{4.6cm}}
\titleformat{\chapter}[block]{}{}{0pt}{\chaptertitleformat}
\titleformat{name=\chapter,numberless}[block]
  {\Huge\bfseries\color{\ChapColor}}{}{0pt}{}
\titlespacing*{\chapter}{0pt}{0pt}{12pt}
\titleformat{\section}{\large\bfseries\color{\ChapColor}}
  {\thesection}{0.7em}{}[\color{\ChapColor}\titlerule]
\titleformat{\subsection}{\normalsize\bfseries\color{DarkText}}
  {\thesubsection}{0.7em}{}
\renewcommand{\chaptermark}[1]{\markboth{\thechapter.\ #1}{}}
\renewcommand{\sectionmark}[1]{\markright{\thesection.\ #1}}
\numberwithin{equation}{chapter}

% Boxes follow the supplied template.
\tcbset{bookbox/.style={
  enhanced,breakable,arc=6pt,boxrule=1pt,
  left=8pt,right=8pt,top=10pt,bottom=7pt,
  fonttitle=\bfseries\small,
  attach boxed title to top left={xshift=8pt,yshift=-3pt},
  boxed title style={arc=4pt,boxrule=0pt}}}
\newenvironment{definition}[1]{%
  \begin{tcolorbox}[bookbox,colback=\ChapColorLight,
    colframe=\ChapColor,colbacktitle=\ChapColor,coltitle=white,
    title={Definition: #1}]}{\end{tcolorbox}}
\newenvironment{example}[1]{%
  \begin{tcolorbox}[bookbox,colback=BoxGray,colframe=RuleGray,
    colbacktitle=MidGray,coltitle=white,
    title={Example: #1}]}{\end{tcolorbox}}
\newenvironment{caution}{%
  \begin{tcolorbox}[bookbox,colback=LightOrange,colframe=Cap2Orange,
    colbacktitle=Cap2Orange,coltitle=white,
    title={Scope of the approximation}]}{\end{tcolorbox}}
\newenvironment{lesson}[1]{%
  \begin{tcolorbox}[bookbox,colback=white,colframe=\ChapColor,
    colbacktitle=\ChapColor,coltitle=white,
    title={#1}]}{\end{tcolorbox}}
\newcommand{\formula}[2][]{%
  \begin{tcolorbox}[enhanced,breakable,
    colback=LightOrange,colframe=Cap2Orange,arc=6pt,boxrule=1.2pt]
    \if\relax\detokenize{#1}\relax\else
      {\small\bfseries\color{Cap2Orange}#1}\par\smallskip
    \fi
    \begin{equation*}#2\end{equation*}
  \end{tcolorbox}}
\setlist[itemize]{label=\textcolor{\ChapColor}{\ding{228}},
  leftmargin=1.4em,itemsep=3pt}
\setlist[enumerate]{leftmargin=1.8em,itemsep=3pt}
\lstset{language=Python,basicstyle=\ttfamily\footnotesize,
  keywordstyle=\color{Cap4Purple},commentstyle=\color{MidGray},
  stringstyle=\color{Cap3Green},backgroundcolor=\color{BoxGray},
  frame=single,rulecolor=\color{RuleGray},breaklines=true,
  columns=fullflexible,keepspaces=true,showstringspaces=false}
\titlecontents{chapter}[0pt]{\vspace{5pt}\bfseries}
  {\contentslabel{1.8em}}{}{\titlerule*[6pt]{.}\contentspage}
\titlecontents{section}[2.5em]{}
  {\contentslabel{2.5em}}{}{\titlerule*[6pt]{.}\contentspage}

\begin{document}

\begin{titlepage}
\begin{tikzpicture}[remember picture,overlay]
  \fill[Cap1Blue]
    (current page.north west) rectangle (current page.south east);
  \fill[Cap2Orange]
    (current page.north west) rectangle
    ($(current page.south west)+(0.65cm,0)$);
  \fill[AccentYellow]
    ($(current page.north west)+(0.65cm,0)$) rectangle
    ($(current page.south west)+(0.95cm,0)$);
  \fill[white,opacity=0.07]
    ($(current page.south east)+(-1cm,3cm)$) circle (6cm);
  \node[anchor=north west,text=white,
    font=\fontsize{42}{48}\selectfont\bfseries,
    align=left,text width=15cm]
    at ($(current page.north west)+(2cm,-3.5cm)$)
    {IONISATION\\PHYSICS};
  \draw[white,line width=1.2pt]
    ($(current page.north west)+(2cm,-7.5cm)$) --
    ($(current page.north east)+(-2cm,-7.5cm)$);
  \node[anchor=north west,text=AccentYellow,
    font=\fontsize{21}{26}\selectfont,align=left,text width=14cm]
    at ($(current page.north west)+(2cm,-8.2cm)$)
    {Atomic processes and radiation\\in the solar atmosphere};
  \node[anchor=north west,text=white,font=\large,
    align=left,text width=14cm]
    at ($(current page.north west)+(2cm,-12cm)$)
    {1.\ Population balance and equilibrium\\[12pt]
     2.\ Ionisation and recombination processes\\[12pt]
     3.\ Radiation, opacity and photoionisation\\[12pt]
     4.\ A rate-based formulation for FIPpy};
  \node[anchor=south west,text=white,font=\small]
    at ($(current page.south west)+(2cm,1.5cm)$)
    {FIPpy \textbullet\ Physical background and proposed formulation};
\end{tikzpicture}
\end{titlepage}

\frontmatter
\chapter*{About these notes}
\addcontentsline{toc}{chapter}{About these notes}
These notes describe the atomic processes that determine the ionisation
fractions used in a FIP fractionation calculation. We begin with the
population balance, then introduce the individual reactions and the
radiation field required for photoionisation. The final chapter combines
these terms into a rate-based formulation for FIPpy.

The rate equations describe a proposed extension of the ionisation
treatment. They should not be read as a statement that all of these
processes are already implemented in FIPpy. In a calculation on a
prescribed atmosphere, the temperature, electron density and hydrogen
populations are inputs. Their consistency and the adopted radiation
field remain part of the physical assumptions.

We use cgs units for densities and atomic coefficients. Binary rate
coefficients have units of $\mathrm{cm^3\,s^{-1}}$, while transition
rates per target particle have units of $\mathrm{s^{-1}}$.

\tableofcontents
\mainmatter

\setChapColor{Cap1Blue}{LightBlue}
\chapter{Population balance and equilibrium}

\section{Neutral and ionised populations}
For an element $X$ with atomic number $Z_X$, let $n_q$ denote the
number density in charge state $q$, where $q=0$ is neutral and
$q=Z_X$ is fully stripped. The total number density and charge-state
fractions are
\begin{equation}
  n_X=\sum_{q=0}^{Z_X}n_q,\qquad
  f_q=\frac{n_q}{n_X},\qquad
  \sum_{q=0}^{Z_X}f_q=1.
\end{equation}

\begin{definition}{Total ionised fraction}
The fraction of an element in all charged states is
\begin{equation}
  \xi_X=\sum_{q=1}^{Z_X}f_q=1-f_0.
  \label{eq:xi}
\end{equation}
If excited levels are resolved, $f_0$ includes all bound levels of the
neutral atom. An atom remains neutral when one of its electrons is
excited to another bound level.
\end{definition}

The distinction between $f_1$ and $\xi_X$ matters when multiply ionised
states are populated. In a neutral--singly-ionised approximation, they
are equal. With additional stages, $\xi_X$ includes their populations
as well.

In the usual low-frequency ponderomotive description, the force acts
on the charged component. The ionisation calculation therefore
determines how much of each element can participate in the
fractionation process \cite{Laming2015}.

\section{Rate coefficients and transition rates}
A binary reaction coefficient describes the effectiveness of a
collision after averaging over the relative velocities of the
particles. Multiplication by the collision-partner density gives the
transition rate per target particle. For electron-impact ionisation,
\begin{equation}
  R_{\mathrm{e,ion},X}=n_e C_X(T_e).
  \label{eq:electron_rate}
\end{equation}
Here $C_X$ has units of $\mathrm{cm^3\,s^{-1}}$, and $R_{\mathrm{e,ion},X}$
has units of $\mathrm{s^{-1}}$. The corresponding number of ionisation
events per unit volume and time is $n_0n_eC_X$, with units of
$\mathrm{cm^{-3}\,s^{-1}}$.

The photoionisation rate $\Gamma_X$ already has units of
$\mathrm{s^{-1}}$. It is calculated from the radiation field and
the atomic cross-section. No additional factor of $n_e$ enters this
rate.

\section{The two-stage balance}
Consider a parcel containing neutral atoms $X^0$ and singly charged
ions $X^+$. Let $R_{\mathrm{ion},X}$ be the total rate per neutral
atom for conversion to $X^+$, and let $R_{\mathrm{rec},X}$ be the
total rate per ion for conversion to $X^0$.

With a conserved number of atoms in the parcel, their ionised
fraction satisfies
\begin{equation}
  \frac{d\xi_X}{dt}
  =(1-\xi_X)R_{\mathrm{ion},X}
  -\xi_X R_{\mathrm{rec},X}.
  \label{eq:two_stage_time}
\end{equation}
Statistical equilibrium sets the left-hand side to zero. Thus,
\begin{equation}
  n_0R_{\mathrm{ion},X}=n_1R_{\mathrm{rec},X},
\end{equation}
and, provided the sum of the rates is positive,
\formula[Two-stage equilibrium]{
  \xi_X^{\mathrm{eq}}
  =\frac{R_{\mathrm{ion},X}}
  {R_{\mathrm{ion},X}+R_{\mathrm{rec},X}}}

Equilibrium fixes the populations by balancing the reactions.
The rates can include electron collisions, photons and charge
exchange in the same calculation.

\section{Collisional equilibrium and LTE}
In a low-density, optically thin collisional model, electron-impact
ionisation is balanced by electron recombination:
\begin{equation}
  R_{\mathrm{ion},X}=n_eC_X,\qquad
  R_{\mathrm{rec},X}=n_e\alpha_X.
\end{equation}
For the two-stage case,
\begin{equation}
  \xi_X^{\mathrm{CIE}}=\frac{C_X}{C_X+\alpha_X}.
\end{equation}
The common factor $n_e$ cancels when the coefficients depend only
on temperature. Density effects can change this result when they
modify the effective coefficients or the level populations.

Local thermodynamic equilibrium (LTE) provides a different closure.
The Saha equation for the first ionisation stage is
\begin{equation}
  \frac{n_1}{n_0}
  =\frac{2}{n_e}
   \left(\frac{2\pi m_ek_BT}{h^2}\right)^{3/2}
   \frac{U_1(T)}{U_0(T)}
   \exp\left(-\frac{\chi_X}{k_BT}\right),
  \label{eq:saha}
\end{equation}
where $\chi_X$ is the first ionisation energy and $U_0,U_1$ are
the internal partition functions. The factor 2 accounts for the
free-electron spin degeneracy in this convention.

Replacing $U_0,U_1$ by ground-state statistical weights is an
additional approximation. Its accuracy depends on the contribution
of excited levels. A rate network containing only electron-impact
ionisation and optically thin radiative recombination does not
automatically recover the Saha limit; the inverse processes required
for detailed balance must also be included.

\section{Equilibrium and time dependence}
For constant two-stage rates, Equation~\eqref{eq:two_stage_time} has
the solution
\begin{equation}
  \xi_X(t)=\xi_X^{\mathrm{eq}}
  +[\xi_X(0)-\xi_X^{\mathrm{eq}}]\exp(-t/\tau_X),
  \qquad
  \tau_X=\frac{1}{R_{\mathrm{ion},X}+R_{\mathrm{rec},X}}.
\end{equation}

Statistical equilibrium is appropriate when the populations relax
on a time scale short compared with changes in the atmosphere and
radiation field. Rapid heating or a flow through a steep atmospheric
gradient can invalidate this assumption. Hydrogen provides a
well-studied example of delayed ionisation and recombination in
the dynamic chromosphere \cite{CarlssonStein2002}.

A non-LTE calculation can still assume statistical equilibrium.
Non-LTE specifies that local thermodynamic relations are
insufficient to determine the populations; time-dependent
ionisation retains their evolution explicitly.

\setChapColor{Cap2Orange}{LightOrange}
\chapter{Atomic processes}

\section{Direct electron-impact ionisation}
A free electron can transfer enough energy to eject a bound
electron:
\begin{equation}
  X^{q+}+e^-\longrightarrow X^{(q+1)+}+2e^-.
\end{equation}
The incident electron must have sufficient energy for the transition.
For thermal electrons, the coefficient $C_{\mathrm{DI},q}(T_e)$
averages the cross-section over their velocity distribution.

At a low electron temperature, relatively few electrons can exceed
a large ionisation threshold. Increasing the temperature changes
the number of electrons able to drive the reaction. The
temperature dependence is already included in the coefficient
returned by fiasco \cite{FiascoDocs}.

\section{Excitation--autoionisation}
An electron collision can first excite an ion into an unstable
state above an ionisation threshold:
\begin{align}
  X^{q+}+e^- &\longrightarrow (X^{q+})^{**}+e^-,\\
  (X^{q+})^{**} &\longrightarrow X^{(q+1)+}+e^-.
\end{align}
The double asterisk denotes the unstable intermediate state.
Its decay ejects an electron, so the net result is ionisation.

Excitation to an ordinary bound level leaves the charge unchanged.
Excitation--autoionisation specifically requires a state that can
decay by electron ejection. Fiasco combines direct ionisation and
this indirect contribution in its total electron-impact coefficient:
\begin{equation}
  C_q=C_{\mathrm{DI},q}+C_{\mathrm{EA},q}.
\end{equation}

\section{Radiative recombination}
A free electron can be captured by an ion while a photon carries
away the excess energy:
\begin{equation}
  X^{(q+1)+}+e^-\longrightarrow X^{q+}+h\nu.
\end{equation}
Capture may populate an excited bound level, followed by further
radiative decays. A total radiative recombination coefficient sums
the relevant capture contributions.

The word \emph{radiative} describes photon emission during capture.
The standard spontaneous recombination coefficient depends on the
electron distribution and atomic structure; it does not require an
incident ionising radiation field. Stimulated contributions require
a separate treatment when they are relevant.

\section{Dielectronic recombination}
During dielectronic capture, a free electron is captured while an
electron already bound to the ion is excited:
\begin{equation}
  X^{(q+1)+}+e^-\longrightarrow (X^{q+})^{**}.
\end{equation}
This intermediate state can either eject an electron or emit a
photon that stabilises the captured electron:
\begin{align}
  (X^{q+})^{**}
    &\longrightarrow X^{(q+1)+}+e^-,
    &&\text{autoionisation},\\
  (X^{q+})^{**}
    &\longrightarrow (X^{q+})^{*}+h\nu,
    &&\text{radiative stabilisation}.
\end{align}
Recombination succeeds through the stabilising branch. The process
involves the incoming electron and an initially bound electron;
it does not describe the capture of two free electrons.

In a resolved calculation, these channels and their branching
probabilities enter explicitly. A total dielectronic coefficient
includes their net contribution. At finite density, collisions can
reionise weakly bound recombined states before they complete their
radiative decays, reducing the effective recombination coefficient
\cite{Nikolic2018}.

For the usual decomposition,
\begin{equation}
  \alpha_q=\alpha_{\mathrm{RR},q}
           +\alpha_{\mathrm{DR},q}.
\end{equation}
Fiasco also uses dedicated total-recombination data for some ions.
Its \texttt{recombination\_rate} property should therefore be used
for the preferred total, even when the components are plotted
separately \cite{FiascoDocs}.

\section{Photoionisation}
A photon can eject an electron when its energy exceeds the
ionisation threshold of the initial state:
\begin{equation}
  X^{q+}+h\nu\longrightarrow X^{(q+1)+}+e^-,
  \qquad h\nu\geq\chi_q.
\end{equation}

The photons can originate outside the local plasma parcel.
Consequently, cool material can be photoionised even when its
thermal electrons rarely reach the required energy. The local
electron temperature and density are insufficient to determine
this contribution.

Photospheric radiation can ionise species with sufficiently low
thresholds. Radiation from the transition region and corona, and
radiation produced within the chromosphere, can also contribute.
Their relative importance depends on the spectrum at the position
of the atom and the cross-section of the transition
\cite{Laming2015,Dufresne2021}.

\section{Charge exchange with hydrogen}
Charge exchange transfers a bound electron between collision
partners. For the first ionisation stage, the reactions are
\begin{align}
  X^0+\mathrm{H}^+&\longrightarrow X^++\mathrm{H}^0,
    &&\text{ionisation of }X,\\
  X^++\mathrm{H}^0&\longrightarrow X^0+\mathrm{H}^+,
    &&\text{recombination of }X.
\end{align}
The electron remains bound to one of the collision partners.
The total charge is conserved, and the reaction itself does not
create a free electron.

\formula[Rates per target particle]{
  \begin{aligned}
    R_{\mathrm{cx,ion},X}
      &=n_{\mathrm{H}^+}k_{\mathrm{cx,ion},X},\\
    R_{\mathrm{cx,rec},X}
      &=n_{\mathrm{H}^0}k_{\mathrm{cx,rec},X}.
  \end{aligned}}

The coefficients depend on the species and the collision conditions.
They describe heavy-particle collisions, so the relevant relative
velocity distribution must be specified when the electron and ion
temperatures differ.

\begin{example}{Oxygen and hydrogen}
The first ionisation energies of oxygen and hydrogen are close.
Charge exchange between them is therefore nearly resonant and can
strongly couple their ionisation states. If these reactions dominate
a two-stage balance,
\begin{equation}
  \frac{n_{\mathrm{O}^+}}{n_{\mathrm{O}^0}}
  \simeq
  \frac{n_{\mathrm{H}^+}}{n_{\mathrm{H}^0}}
  \frac{k_{\mathrm{cx,ion},\mathrm O}}
       {k_{\mathrm{cx,rec},\mathrm O}}.
\end{equation}
The coupling does not require the O and H ionised fractions to be
exactly equal.
\end{example}

Charge exchange can also affect other elements. Laming's later
treatment revises the adopted rates for sulphur, sodium and iron,
illustrating the importance of the species-specific data
\cite{Laming2021}. Elastic collisions that exchange momentum and
charge-transfer reactions should retain separate coefficients.

\setChapColor{Cap3Green}{LightGreen}
\chapter{Radiation and opacity}

\section{Specific intensity and mean intensity}
The specific intensity $I_\nu(s,\Omega)$ describes radiation at
frequency $\nu$, position $s$ and propagation direction $\Omega$.
In cgs units,
\begin{equation}
  [I_\nu]=
  \mathrm{erg\,s^{-1}\,cm^{-2}\,sr^{-1}\,Hz^{-1}}.
\end{equation}
The mean intensity averages over directions:
\begin{equation}
  J_\nu(s)=\frac{1}{4\pi}\int I_\nu(s,\Omega)\,d\Omega.
  \label{eq:mean_intensity}
\end{equation}

Photons arriving from any direction can ionise an atom. This is why
the photoionisation rate uses the angular integral of intensity.
The flux through a chosen surface instead contains the directional
factor $\mu$, the cosine relative to its normal:
\begin{equation}
  F_\nu=\int I_\nu(\Omega)\mu\,d\Omega.
\end{equation}
An isotropic field can have nonzero $J_\nu$ and zero net flux.
A measured irradiance therefore requires a geometrical prescription
before it can be used as a local mean intensity.

\section{The photoionisation integral}
Let $\sigma_X(\nu)$ be the photoionisation cross-section of the
initial atomic state. Dividing the radiation intensity by $h\nu$
converts energy into photon number. The rate per atom is
\begin{equation}
  \Gamma_X(s)=4\pi
    \int_{\nu_{\mathrm{th},X}}^\infty
    \frac{J_\nu(s)}{h\nu}\sigma_X(\nu)\,d\nu,
  \qquad h\nu_{\mathrm{th},X}=\chi_X.
  \label{eq:gamma_nu}
\end{equation}
The result has units of $\mathrm{s^{-1}}$. The photoionisation
events per unit volume and time are $n_0\Gamma_X$ for a neutral
initial population.

For a spectrum tabulated in wavelength, the equivalent expression is
\formula[Photoionisation on a wavelength grid]{
  \Gamma_X(s)=4\pi
    \int_0^{\lambda_{\mathrm{th},X}}
    J_\lambda(s)\frac{\lambda}{hc}
    \sigma_X(\lambda)\,d\lambda,
  \qquad
  \lambda_{\mathrm{th},X}=\frac{hc}{\chi_X}.}
Here $J_\lambda$ is a spectral intensity per unit wavelength.
The spectral-density conversion satisfies
$J_\lambda\,|d\lambda|=J_\nu\,|d\nu|$; replacing the axis without
converting the spectral density changes the integral.

Ground-state cross-sections can be obtained from the fits of
Verner et al.\ \cite{Verner1996} within their stated coverage and
energy ranges. An excited initial level has its own threshold
and cross-section, which must be included in a level-resolved model.

\begin{example}{Photons above the threshold}
For a narrow incident line at frequency $\nu_L$, let
$\mathcal{J}_L=\int_{\mathrm{line}}J_\nu\,d\nu$. If the cross-section
varies little across the line and $\nu_L>\nu_{\mathrm{th},X}$,
\begin{equation}
  \Gamma_{X,L}\simeq
  \frac{4\pi\mathcal{J}_L}{h\nu_L}\sigma_X(\nu_L).
\end{equation}
If the line lies below the initial state's threshold, its photons
cannot photoionise that state. A bright line therefore need not
contribute to every element.
\end{example}

\section{Opacity and the propagation of photons}
The extinction coefficient $\kappa_\nu$ has units of inverse
length. Optical depth along a radiation path $\ell$ is
\begin{equation}
  \tau_\nu(\ell_0,\ell)=
  \int_{\ell_0}^{\ell}\kappa_\nu(\ell')\,d\ell'.
\end{equation}
For a bound-free absorber $a$,
\begin{equation}
  \kappa_{\nu,a}^{\mathrm{bf}}=n_a\sigma_a(\nu).
\end{equation}
Several populations can contribute to the total opacity.

If a ray undergoes pure absorption and there is no emission along
its path, its intensity is
\begin{equation}
  I_\nu(\ell)=I_{\nu,0}\exp[-\tau_\nu(\ell_0,\ell)].
  \label{eq:attenuation}
\end{equation}
The incident intensity $I_{\nu,0}$ is an independent input.
Opacity specifies attenuation; it does not specify the source
brightness.

With emission included, the transfer equation is
\begin{equation}
  \frac{dI_\nu}{d\ell}=-\kappa_\nu I_\nu+j_\nu,
  \label{eq:transfer}
\end{equation}
where $j_\nu$ is the emissivity. Scattering contributes to extinction
and to radiation entering the ray from other directions. In that
case the radiation field must be found with the corresponding
scattering treatment.

\begin{lesson}{From a radiation source to an ionisation rate}
A photoionisation calculation requires the following quantities:
\begin{enumerate}
  \item Boundary illumination and any internal emissivity.
  \item Absorption and scattering along the radiation paths.
  \item The local angular mean $J_\nu$ obtained from the resulting intensities.
  \item The photoionisation cross-section in
    Equation~\eqref{eq:gamma_nu}.
\end{enumerate}
The same cross-section can contribute to opacity through
$n_X\sigma_X$ and to the rate per atom through the integral over
$J_\nu\sigma_X/(h\nu)$.
\end{lesson}

\section{Hydrogen attenuation and trapped radiation}
A first approximation may retain only ground-state neutral
hydrogen continuum absorption:
\begin{equation}
  \tau_\nu=\sigma_{\mathrm{H\,I}}(\nu)
  \int n_{\mathrm{H}^0}\,d\ell.
\end{equation}
This continuum cross-section vanishes below the H ground-state
ionisation threshold of approximately $13.6\,\mathrm{eV}$.
Metal continua and spectral lines can still provide opacity
at those photon energies.

In particular, the ground-state H photoionisation cross-section
does not describe resonant scattering in Ly$\alpha$. An incident
Ly$\alpha$ component and locally produced, trapped Ly$\alpha$
require a line-radiation prescription. Their effects on
photoionisation depend on the thresholds of the absorbing states
\cite{Laming2015}.

Radiation below the hydrogen continuum threshold is relevant to
low-FIP species. Improved C, Si and S atomic models can change the
UV opacity and the resulting radiation field
\cite{Dufresne2025}.

\section{Boundary illumination and geometry}
For a uniformly bright photospheric surface, immediately above
that surface the upward hemisphere has
\begin{equation}
  F_{\nu,\mathrm{surf}}=\pi I_{\nu,\mathrm{surf}},
  \qquad J_{\nu,\mathrm{ph}}=\frac{I_{\nu,\mathrm{surf}}}{2}.
\end{equation}
Under the same approximation, an irradiance measured at distance
$d$ gives
\begin{equation}
  J_{\nu,\mathrm{ph}}
  =\frac{F_{\nu,d}}{2\pi}
   \left(\frac{d}{R_\odot}\right)^2.
\end{equation}
The factor $1/2$ averages a uniform intensity over an illuminated
hemisphere. The assumptions include the adopted surface brightness
and the near-surface geometry.

Downward illumination from the corona requires its own source
geometry. A disk-integrated outward irradiance does not uniquely
specify the radiation received from the overlying corona.

The loop coordinate $s$ follows the magnetic field. Radiation may
travel across that field, so an opacity integral along $s$ represents
a specified path approximation. Treating
$J_\nu=J_{\nu,0}\exp(-\tau_\nu)$ as an attenuation law for the mean
intensity similarly requires an effective angular prescription.
A transfer solution evaluates the ray-dependent optical depths
before averaging over directions.

\setChapColor{Cap4Purple}{LightPurple}
\chapter{A formulation for FIPpy}

\section{Combining the two-stage rates}
For the neutral--singly-ionised balance, define
\begin{align}
  R_{\mathrm{ion},X}
    &=n_eC_X+\Gamma_X
      +n_{\mathrm{H}^+}k_{\mathrm{cx,ion},X},
      \label{eq:total_ion}\\
  R_{\mathrm{rec},X}
    &=n_e\alpha_X
      +n_{\mathrm{H}^0}k_{\mathrm{cx,rec},X}.
      \label{eq:total_rec}
\end{align}
Here $C_X$ describes ionisation out of $X^0$, while $\alpha_X$
describes recombination out of $X^+$ into $X^0$. Each term on the
right-hand sides has units of $\mathrm{s^{-1}}$.

The statistical-equilibrium fraction is then
\begin{equation}
  \xi_X^{\mathrm{eq}}=
  \frac{n_eC_X+\Gamma_X+
        n_{\mathrm{H}^+}k_{\mathrm{cx,ion},X}}
       {R_{\mathrm{ion},X}+R_{\mathrm{rec},X}}.
  \label{eq:full_two_stage}
\end{equation}
Additional collision partners, including helium, can be introduced
through their own populations and reaction coefficients.

The relative importance of the ionisation channels can be
diagnosed by comparing $n_eC_X$, $\Gamma_X$ and
$n_{\mathrm{H}^+}k_{\mathrm{cx,ion},X}$ at each location.
The final FIP sensitivity also depends on where these reactions
change $\xi_X$ relative to the ponderomotive acceleration and
collision rates.

\section{What fiasco provides}
Fiasco accesses CHIANTI atomic data and calculates electron-impact
ionisation and electron recombination coefficients
\cite{FiascoDocs,Dufresne2024}. For a neutral--singly-ionised
calculation, the required quantities belong to adjacent ion objects.

\begin{example}{Selecting the oxygen coefficients}
The following code illustrates the stage selection. It requires
fiasco and an available CHIANTI database.
\begin{lstlisting}
import astropy.units as u
import numpy as np
from fiasco import Ion

# Illustrative grid; check the atomic fits' valid ranges.
T = np.logspace(3.5, 6, 200) * u.K
oxygen_neutral = Ion("O I", T)
oxygen_ion = Ion("O II", T)

# O I -> O II
C_O = oxygen_neutral.ionization_rate

# O II -> O I
alpha_O = oxygen_ion.recombination_rate
\end{lstlisting}
Both quantities are coefficients with units of
$\mathrm{cm^3\,s^{-1}}$. Multiplication by the local electron
density gives their contributions to
Equations~\eqref{eq:total_ion} and \eqref{eq:total_rec}.
\end{example}

For an Fe XVI object, the ionisation property describes
Fe XVI to Fe XVII, and the recombination property describes
Fe XVI to Fe XV. Plotting both is useful for examining the
loss channels of Fe XVI, but those two curves do not form
the forward and reverse rates of a single adjacent-stage pair.

The usual temperature-dependent properties do not by themselves
supply the local radiation field, photoionisation or H charge
exchange. Support for a CHIANTI database version also does not
establish that every advanced density-dependent process is
included by those properties. CHIANTI v11 provides advanced
models for selected elements; the applicable dataset and
software treatment must be identified \cite{Dufresne2024}.

\section{Several charge states}
Let $I_q$ be the total rate for $q\rightarrow q+1$ and $A_q$
the total rate for $q\rightarrow q-1$. For an adjacent-stage
network following a parcel,
\begin{equation}
  \frac{df_q}{dt}
  =I_{q-1}f_{q-1}+A_{q+1}f_{q+1}
   -(I_q+A_q)f_q.
  \label{eq:multistage}
\end{equation}
Terms outside the allowed charge range are absent, with
$A_0=0$ and $I_{Z_X}=0$. This form assumes that the represented
reactions change the charge by one.

For a prescribed background, the same equations can be written
\begin{equation}
  \frac{d\boldsymbol f}{dt}=\boldsymbol R\boldsymbol f.
\end{equation}
Using column-vector populations, an off-diagonal entry $R_{ij}$
is the rate from state $j$ to state $i$. Each diagonal entry
is minus the sum of the rates leaving that state, so the
columns of $\boldsymbol R$ sum to zero.

Statistical equilibrium requires
\begin{equation}
  \boldsymbol R\boldsymbol f=0,\qquad \sum_q f_q=1.
\end{equation}
A level-resolved model uses the same construction with individual
atomic levels as states. The neutral fraction is then the sum
over all neutral levels.

\section{Prescribed atmospheres and time dependence}
On a prescribed HYDRAD atmosphere, $T_e$, $n_e$ and the hydrogen
populations are supplied to the metal rate solver. This treats
the metals as trace species for the background calculation.
The resulting populations do not feed back on the saved
temperature or electron density.

A coupled calculation instead has to update the electron budget,
energy balance and radiation consistently. Changing $n_e$
while holding the rest of a saved atmosphere fixed introduces a
different physical closure.

For evolving atmospheres, Equation~\eqref{eq:multistage} retains
the population history. Transport with a common flow speed $u$
in a tube of cross-sectional area $\mathcal A$ can be written
\begin{equation}
  \frac{\partial(\mathcal A n_q)}{\partial t}
  +\frac{\partial(\mathcal A u n_q)}{\partial s}
  =\mathcal A\sum_jR_{qj}n_j.
\end{equation}
Initial populations, boundary inflows and sufficiently resolved
atmospheric histories are required. Differential ion--neutral
transport and the formation of an abundance anomaly require
their own transport treatment.

\section{Assumptions to record with each calculation}
\begin{caution}
Equations~\eqref{eq:total_ion}--\eqref{eq:full_two_stage} describe
a two-stage, statistical-equilibrium model on a specified
background. They omit higher charge stages unless the network
is extended. A prescribed attenuated spectrum approximates the
radiation field; it does not solve the full non-LTE transfer
problem. Time dependence, density effects and excited-level
populations enter only when their corresponding rates and
equations are included.
\end{caution}

For each run, record the atomic-data version, reaction coverage
and validity ranges. State the source of the H and He populations,
the treatment of electron density, and the boundary radiation.
Avoid applying a density-suppression correction to recombination
coefficients that already contain that correction.

The calculation should conserve the total population, return
nonnegative fractions and satisfy the equilibrium residual when
equilibrium is requested. Numerical checks should also include
spectral-grid convergence and, for time-dependent runs, time-step
and atmospheric-cadence convergence.

A useful first implementation is the collisional rate matrix,
followed by O--H charge exchange and photoionisation. Other
elements then require their own verified reaction data.
The resulting neutral fractions and FIP predictions can be
compared with the existing Saha--CHIANTI prescription on the
same atmosphere.

\backmatter
\setChapColor{Cap1Blue}{LightBlue}
\begin{thebibliography}{99}
\addcontentsline{toc}{chapter}{References}
\bibitem{Laming2015}
Laming, J. M. (2015).
\textit{The FIP and Inverse FIP Effects in Solar and Stellar Coronae}.
Living Reviews in Solar Physics, 12, 2.
\url{https://doi.org/10.1007/lrsp-2015-2}.

\bibitem{Laming2021}
Laming, J. M. (2021).
\textit{The FIP and Inverse-FIP Effects in Solar Flares}.
\url{https://doi.org/10.3847/1538-4357/abd9c3}.

\bibitem{CarlssonStein2002}
Carlsson, M., and Stein, R. F. (2002).
\textit{Dynamic Hydrogen Ionization}.
\url{https://arxiv.org/abs/astro-ph/0202313}.

\bibitem{Verner1996}
Verner, D. A., Ferland, G. J., Korista, K. T., and Yakovlev, D. G.
(1996).
\textit{Atomic Data for Astrophysics. II. New Analytic Fits for
Photoionization Cross Sections of Atoms and Ions}.
\url{https://arxiv.org/abs/astro-ph/9601009}.

\bibitem{Nikolic2018}
Nikoli\'c, D., et al. (2018).
\textit{Suppression of Dielectronic Recombination Due to Finite
Density Effects II: Analytical Refinement and Application to
Density-dependent Ionization Balances and AGN Broad-line Emission}.
\url{https://arxiv.org/abs/1808.07365}.

\bibitem{Dufresne2021}
Dufresne, R. P., Del Zanna, G., and Badnell, N. R. (2021).
\textit{The influence of photo-induced processes and charge transfer
on carbon and oxygen in the lower solar atmosphere}.
\url{https://arxiv.org/abs/2102.09278}.

\bibitem{Dufresne2024}
Dufresne, R. P., et al. (2024).
\textit{CHIANTI: An Atomic Database for Emission Lines.
Paper XVIII. Version 11, Advanced Ionization Equilibrium Models:
Density and Charge Transfer Effects}.
\url{https://arxiv.org/abs/2403.16922}.

\bibitem{Dufresne2025}
Dufresne, R. P., Del Zanna, G., and Osborne, C. M. J. (2025).
\textit{Improving the atomic modelling for solar UV radiative
transfer calculations}.
\url{https://doi.org/10.1093/mnras/staf1330}.

\bibitem{FiascoDocs}
\textit{Fiasco documentation: Ion API and ionisation/recombination rates}.
\url{https://fiasco.readthedocs.io/en/stable/api/fiasco.Ion.html}.
\end{thebibliography}
\end{document}

